\let\cleardoublepage\clearpage
\addcontentsline{toc}{chapter}{Resumen}
\renewcommand{\abstractname}{Resumen}


\thispagestyle{plain}
\begin{center}
    \textbf{MODELADO DE SERIES TEMPORALES ECONÓMICAS: DE LA TASA DE CAMBIO RELATIVA A LOS MODELOS DE TRANSICIÓN DE RÉGIMEN ESTOCÁSTICAMENTE ESTRUCTURADOS}

        \vspace{0.9cm}
    \textbf{Resumen}
\end{center}

El análisis de regímenes en series temporales económicas y financieras presenta una disyuntiva persistente entre modelos deterministas, que carecen de flexibilidad estocástica, y modelos probabilísticos de difícil interpretación.  Esta investigación propone un marco unificado para cerrar esa brecha, articulado en tres contribuciones progresivas.

La primera contribución es la formalización de la familia de operadores de Tasa de Cambio Relativa Óptima Conmutada (TCROC) y sus variantes paramétricas (TCROCM, ETCROC, ETCROCM), con demostraciones rigurosas de existencia, unicidad, consistencia asintótica y estabilidad ante perturbaciones.  Estos operadores transforman una serie temporal continua en una secuencia de regímenes económicos con significado interpretable.

La segunda contribución es la integración del operador TCROC con cadenas de Markov y la estimación de la matriz de transición mediante un algoritmo de Mínimos Cuadrados No Negativos (NNLS).  A diferencia del enfoque de máxima verosimilitud utilizado en los modelos Markov-Switching (MS-AR), la formulación NNLS es convexa, garantiza convergencia global, y asegura por construcción la validez estocástica de la matriz estimada.

La tercera contribución es la extensión del marco lineal mediante Computación de Reservorio Estocásticamente Estructurada (SSRC).  Se demuestra formalmente que el modelo TCROC-Markov es un caso particular de los SSRC (Teorema de Inclusión), proporcionando un camino teórico hacia la captura de dinámicas no lineales.  La validación empírica de esta extensión muestra mejoras del 20--23\% en precisión de pronóstico respecto al modelo lineal.

La validación empírica integral se realiza sobre series semanales de cuatro combustibles hondureños (Regular, Súper, Diésel y Kerosene) durante 2017--2025, incluyendo optimización exhaustiva de hiperparámetros, análisis espectral y comparación con benchmarks (cuantiles fijos y MS-AR).  Los resultados confirman la superioridad estadísticamente significativa del modelo TCROC-Markov y revelan su doble utilidad: alta precisión predictiva con datos extensos y captura de estructura dinámica con datos limitados.

\textbf{Palabras clave:} TCROC, series temporales, cadenas de Markov, NNLS, cambio de régimen, computación de reservorio, SSRC, combustibles, Honduras.
